% from ega4-16-fr.tex
\oldpage[IV-4]{21}
Ce lemme étant démontré, il reste à voir que l'homomorphisme (16.4.8.1) est
bien l'image par $f^*$ de l'homomorphisme (16.4.8.4) correspondant à la suite
exacte (16.4.8.2); on constate aisément que cela résulte de la définition de
$u$ (16.4.8.2) et de celle de $\delta$ (16.4.7.1), compte tenu de la définition
de la structure de $\mathcal{O}_X$-Algèbre sur $\mathscr{P}_{X/S}^n$ et de
celle de $d_{X/S}^n$ (16.3.5 et 16.3.6).

En particulier:

\begin{corollary}[16.4.9]
\label{IV.16.4.9}
Sous les conditions de (16.4.7), on a un isomorphisme canonique
\[
  \operatorname{gr}_1(\delta):f^*(\mathscr{E})
  \xrightarrow{\sim}\Omega_{X/S}^1.
  \tag{16.4.9.1}
  \label{IV.16.4.9.1}
\]
\end{corollary}

\begin{corollary}[16.4.10]
\label{IV.16.4.10}
Si $S=\operatorname{Spec}(A)$, $\mathscr{E}=\mathcal{O}_S^m$, de sorte que
\[
  X=\operatorname{Spec}(A[T_1,\ldots,T_m]),
\]
$\mathscr{P}_{X/S}^n$ s'identifie canoniquement à la
$\mathcal{O}_X$-Algèbre correspondant à la $A[T_1,\ldots,T_m]$-algèbre quotient
\[
  A[T_1,\ldots,T_m,U_1,\ldots,U_m]/\mathfrak{K}^{n+1},
\]
où les $U_i$ ($1\leq i\leq m$) sont $m$ nouvelles indéterminées et
$\mathfrak{K}$ est l'idéal engendré par $U_1,\ldots,U_m$.
\end{corollary}

On retrouve ainsi en particulier la structure de $\Omega_{X/S}^1$ dans ce cas
(0, 20.5.13).

Notons en outre que $d_{X/S}^n$ fait alors correspondre à un polynôme
$F(T_1,\ldots,T_m)$ la classe mod.\ $\mathfrak{K}^{n+1}$ de
$F(T_1+U_1,\ldots,T_m+U_m)$, comme il résulte de la définition (16.4.7.1).

\begin{proposition}[16.4.11]
\label{IV.16.4.11}
Soient $f:X\to S$ un morphisme, $g:S\to X$ une $S$-section de $X$,
$S_g^{(n)}$ le $n$-ème voisinage infinitésimal de $S$ pour l'immersion $g$
(16.1.2). Il existe alors un isomorphisme et un seul de
$\mathcal{O}_S$-Algèbres
\[
  \varpi_n:g^*(\mathscr{P}_{X/S}^n)\longrightarrow
  \mathcal{O}_{S_g^{(n)}}
  \tag{16.4.11.1}
  \label{IV.16.4.11.1}
\]
(pour la structure de $\mathcal{O}_S$-Algèbre sur
$\mathcal{O}_{S_g^{(n)}}$ définie par $f$ (16.1.7)), rendant commutatif le
diagramme
\[
  \xymatrix{
    \mathcal{O}_S=g^*(\mathcal{O}_X)
      \ar[rr]^-{\lambda_n}\ar[dr]_-{g^*(d_{X/S}^n)}
      &&\mathcal{O}_{S_g^{(n)}}\\
    &g^*(\mathscr{P}_{X/S}^n)\ar[ur]_-{\varpi_n}
  }
  \tag{16.4.11.2}
  \label{IV.16.4.11.2}
\]
(où $\lambda_n$ est l'homomorphisme structural).
\end{proposition}

En vertu de (I, 5.3.7), où l'on remplace $X$, $Y$, $S$ par $S$, $X$, $S$
respectivement et $f$ par $g$, les diagrammes
\[
  \xymatrix@=3pc{
    S\ar[r]^-g\ar[d]_-g&X\ar[d]^-{\Delta_f}
      &S\ar[r]^-g\ar[d]_-g&X\ar[d]^-{\Delta_f}\\
    X\ar[r]_{(g\circ f,1_X)_S}&X\times_S X
      &X\ar[r]_{(1_X,g\circ f)_S}&X\times_S X
  }
  \tag{16.4.11.3}
  \label{IV.16.4.11.3}
\]
